\documentclass[11pt]{article}
\usepackage{amsmath, amssymb, amsthm}
\usepackage[margin=1in]{geometry}
\usepackage{algorithm, algpseudocode}

\newcommand{\Gam}{\mathrm{Gam}}
\newcommand{\Poi}{\mathrm{Poi}}
\newcommand{\Multi}{\mathrm{Multi}}
\newcommand{\KL}{\mathrm{KL}}
\newcommand{\E}{\mathbb{E}}

\theoremstyle{definition}
\newtheorem{remark}{Remark}

\title{Alternating Hard-Motif MoS-Poisson-SuSiE}
\author{}
\date{}

\begin{document}
\maketitle

\section{Purpose}

This note describes the current implemented mixture-of-submanifolds
Poisson-SuSiE algorithm. It is an intermediate model between the
MF-Poisson-SuSiE model and the final soft submanifold model. The current version
uses hard submanifold motifs and alternates between:
\begin{enumerate}
  \item fitting observation-level mixture responsibilities,
  \item updating the shared factor matrix $F$, and
  \item updating hard submanifold motifs.
\end{enumerate}

The hard motifs make the algorithm easy to debug and provide a concrete test of
whether a submanifold layer can improve the tree-like sparsity scenarios.

\section{Model}

We observe a count matrix $Y \in \mathbb{N}^{N \times M}$. Let
$F \in \mathbb{R}_+^{K \times M}$ be a row-normalized factor matrix:
\[
  F_{km} \geq 0, \qquad \sum_{m=1}^{M} F_{km} = 1, \quad k \in [K].
\]

There are $S$ submanifolds. Each submanifold $s \in [S]$ has $D$ hard motif
slots,
\[
  a_{sd} \in [K], \qquad d \in [D],
\]
where $a_{sd}$ is the factor used by slot $d$ of submanifold $s$. Repeated
entries are allowed. Thus a submanifold with effective dimension one may be
represented as $(k,k,\ldots,k)$.

For observation $i$, let $z_i \in [S]$ denote its submanifold assignment. The
current hard-motif generative model is:
\begin{align}
  z_i &\sim \Multi(1,\pi), \\
  \beta_{isd} \mid z_i=s &\sim \Gam(\alpha^0_{sd}, \lambda^0_{sd}),
    \qquad d \in [D], \\
  y_{im} \mid z_i=s, \beta_{is1:D}, F, a
    &\sim \Poi\left(\sum_{d=1}^{D} \beta_{isd} F_{a_{sd},m}\right),
    \qquad m \in [M].
\end{align}

This can be viewed as a mixture model over sparse Poisson factorization
submanifolds. Conditional on $z_i=s$, the observation can only use factors
listed in the hard motif $a_{s1:D}$.

\section{Variational family}

The current implementation uses:
\[
  q(z,\beta) =
  \prod_{i=1}^{N} q(z_i)
  \prod_{s=1}^{S}\prod_{d=1}^{D}
  q(\beta_{isd} \mid z_i=s),
\]
where
\[
  q(z_i=s) = \omega_{is}, \qquad
  q(\beta_{isd} \mid z_i=s) = \Gam(\alpha_{isd}, \lambda_{isd}).
\]

For each candidate submanifold $s$, auxiliary variables
$\xi_{isdm} \geq 0$ are introduced with
\[
  \sum_{d=1}^{D} \xi_{isdm} = 1, \qquad i \in [N],\; s \in [S],\; m \in [M].
\]
They allocate count $y_{im}$ across the $D$ motif slots of submanifold $s$.

\section{Conditional component updates}

Fix $F$, the hard motif matrix $a$, and prior parameters
$(\alpha^0,\lambda^0)$. For each submanifold $s$, the conditional
Poisson-SuSiE update has the same form as the fixed-motif single-observation
Poisson update.

\paragraph{Auxiliary allocation.}
For each $i,s,m,d$,
\[
  \xi_{isdm}
  \leftarrow
  \frac{
    \exp\{\E_q[\log \beta_{isd}] + \log F_{a_{sd},m}\}
  }{
    \sum_{d'=1}^{D}
    \exp\{\E_q[\log \beta_{isd'}] + \log F_{a_{sd'},m}\}
  },
\]
where
\[
  \E_q[\log \beta_{isd}] = \psi(\alpha_{isd}) - \log \lambda_{isd}.
\]

\paragraph{Gamma posterior.}
For each $i,s,d$,
\begin{align}
  \alpha_{isd}
    &\leftarrow \sum_{m=1}^{M} \xi_{isdm} y_{im} + \alpha^0_{sd}, \\
  \lambda_{isd}
    &\leftarrow \sum_{m=1}^{M} F_{a_{sd},m} + \lambda^0_{sd}.
\end{align}
Since $F$ is row-normalized, $\sum_m F_{a_{sd},m}=1$ in the current
implementation.

\paragraph{Component lower bound.}
Let $\mathcal{L}_{is}$ denote the conditional variational lower bound for
observation $i$ under submanifold $s$. Up to constants not depending on $s$,
the implemented quantity is
\begin{align}
  \mathcal{L}_{is}
  &=
  \sum_{d,m} \xi_{isdm} y_{im}
    \left\{\E_q[\log \beta_{isd}] + \log F_{a_{sd},m}\right\}
  - \sum_{d,m} y_{im}\xi_{isdm}\log \xi_{isdm} \nonumber \\
  &\quad
  - \sum_d \E_q[\beta_{isd}] \sum_m F_{a_{sd},m}
  - \sum_d \KL\left\{
      \Gam(\alpha_{isd},\lambda_{isd})
      \,\|\, \Gam(\alpha^0_{sd},\lambda^0_{sd})
    \right\}.
\end{align}

\section{Mixture responsibility update}

Given component lower bounds $\mathcal{L}_{is}$, the observation-level
responsibility update is
\[
  \omega_{is}
  \leftarrow
  \frac{\pi_s \exp(\mathcal{L}_{is})}
       {\sum_{s'=1}^{S} \pi_{s'} \exp(\mathcal{L}_{is'})}.
\]
The mixing weights are updated by
\[
  \pi_s \leftarrow \frac{1}{N}\sum_{i=1}^{N}\omega_{is}.
\]

\section{Empirical-Bayes prior update}

The current implementation optionally updates the motif-slot Gamma prior
parameters. For each submanifold $s$ and slot $d$, compute responsibility-
weighted posterior moments:
\begin{align}
  \bar{\beta}_{sd}
    &= \frac{\sum_i \omega_{is}\,\alpha_{isd}/\lambda_{isd}}
            {\sum_i \omega_{is}}, \\
  \overline{\log \beta}_{sd}
    &= \frac{\sum_i \omega_{is}\,
        \{\psi(\alpha_{isd})-\log\lambda_{isd}\}}
            {\sum_i \omega_{is}}.
\end{align}
The Gamma shape solves
\[
  \log \alpha^0_{sd} - \psi(\alpha^0_{sd})
  =
  \log \bar{\beta}_{sd} - \overline{\log \beta}_{sd},
\]
using Newton's method, and then
\[
  \lambda^0_{sd} \leftarrow \frac{\alpha^0_{sd}}{\bar{\beta}_{sd}}.
\]

\section{Shared factor update}

Given current responsibilities $\omega$, auxiliary allocations $\xi$, and hard
motifs $a$, the row-normalized $F$ update collects responsibility-weighted
expected counts:
\[
  A_{km}
  =
  \sum_{i=1}^{N}\sum_{s=1}^{S}\sum_{d=1}^{D}
  \omega_{is}\,\mathbf{1}\{a_{sd}=k\}\,\xi_{isdm}\,y_{im}.
\]
Then
\[
  F^{\mathrm{new}}_{km}
  \leftarrow
  \frac{A_{km}+\epsilon}
       {\sum_{m'=1}^{M} (A_{km'}+\epsilon)}.
\]
The implementation supports a damped update
\[
  F_{k\cdot}
  \leftarrow
  (1-\eta_t)F_{k\cdot} + \eta_t F^{\mathrm{new}}_{k\cdot},
  \qquad
  F_{k\cdot} \leftarrow
  \frac{F_{k\cdot}}{\sum_m F_{km}},
\]
where $\eta_t$ is ramped from an initial value to $1$.

\section{Hard motif update}

The current hard motif update is intentionally simple. It first aggregates
counts by submanifold responsibility:
\[
  Y^{(S)}_{sm} = \sum_{i=1}^{N} \omega_{is} y_{im}.
\]
For fixed $F$, it then fits nonnegative submanifold loadings
$H \in \mathbb{R}_+^{S \times K}$ by KL-NMF multiplicative updates:
\[
  H_{sk}
  \leftarrow
  H_{sk}
  \frac{\sum_m (Y^{(S)}_{sm}/(HF)_{sm})F_{km}}
       {\sum_m F_{km}}.
\]
The initialization of $H$ uses the current hard motif counts, so the update is
local around the current motif state.

After fitting $H$, define normalized submanifold loading shares
\[
  r_{sk} = \frac{H_{sk}}{\sum_{\ell=1}^{K}H_{s\ell}}.
\]
For each submanifold $s$, the updated hard motif is constructed by selecting
factors with $r_{sk}$ above a threshold $\tau$, keeping at most $D$ factors, and
padding with repeats if fewer than $D$ factors are selected:
\[
  a_{s1:D}
  \leftarrow
  \mathrm{Pad}_D\left(
    \{k: r_{sk}\geq \tau\}
  \right).
\]
This allows the fitted effective submanifold dimension to be smaller than the
nominal $D$.

\begin{remark}[Why this motif update is not the final model]
This update learns hard motifs, not posterior motif probabilities. It is useful
as an intermediate coordinate step, but it does not yet quantify uncertainty
over submanifold-level supports. The final model should replace the hard motif
matrix $a$ with soft submanifold-level posterior probabilities over factors.
\end{remark}

\section{Algorithm}

\begin{algorithm}[H]
\caption{Alternating hard-motif MoS-Poisson-SuSiE}
\small
\begin{algorithmic}[1]
\State Initialize $F$ by running MF-Poisson-SuSiE with learned $F$.
\State Compute posterior loading scores from the MF-Poisson-SuSiE fit:
\[
  \widehat{L}_{ik}
  =
  \sum_{d=1}^{D}
  q(\gamma_{id}=k)\,
  \E_q[\beta_{id}\mid \gamma_{id}=k].
\]
\State Cluster row-normalized $\widehat{L}$ into $S$ groups.
\State Initialize hard motifs $a_{s1:D}$ from the cluster centers, allowing
  effective dimension smaller than $D$ by thresholding and padding with repeats.
\State Initialize $\pi_s=1/S$ and $\xi_{isdm}=1/D$.
\For{$t=1,\ldots,T$}
  \State \textbf{Conditional component update:}
  for every submanifold $s$, update $\xi_{isdm}$, $\alpha_{isd}$, and
  $\lambda_{isd}$ using the current $F$ and $a$.
  \State \textbf{Mixture update:}
  compute component lower bounds $\mathcal{L}_{is}$ and update
  responsibilities $\omega_{is}$.
  \State Update mixing proportions
  $\pi_s \leftarrow N^{-1}\sum_i \omega_{is}$.
  \State \textbf{Prior update} (optional):
  update $(\alpha^0_{sd},\lambda^0_{sd})$ using the responsibility-weighted
  Gamma moment equations.
  \State \textbf{Factor update} (optional):
  compute $A_{km}$ and update row-normalized $F$ using a damped step.
  \State \textbf{Motif update} (optional):
  aggregate $Y^{(S)}_{sm}=\sum_i\omega_{is}y_{im}$, fit submanifold loading
  shares $r_{sk}$ with fixed-$F$ KL-NMF updates, and update hard motifs
  $a_{s1:D}$ by thresholding and padding.
  \State Record the mixture lower bound
  $\sum_i \log\{\sum_s \pi_s\exp(\mathcal{L}_{is})\}$.
\EndFor
\State Refit the conditional components and responsibilities one final time
using the returned $F$ and motifs $a$.
\end{algorithmic}
\end{algorithm}

\section{Current interpretation}

The current algorithm should be interpreted as a pragmatic hard-coordinate
approximation. It answers the question:

\begin{quote}
Can a submanifold mixture layer, initialized from MF-Poisson-SuSiE, improve
recovery of repeated sparsity patterns and tree-like support structures?
\end{quote}

The present implementation does learn $F$ and hard motifs alternately, but only
after an MF-Poisson-SuSiE initialization. The motif update is hard and local,
and therefore the ELBO recorded before the $F$ and motif updates is best treated
as a diagnostic rather than as a guaranteed monotone objective. The next
theoretical step is to derive and implement the soft motif posterior version of
the model.

\end{document}
